%!TeX root=../treasuretop.tex
\ArtChapter[How the Ship was Abandoned]{chap16head}

   \lettrine[lines=4,image=true]{chap16initI}{t} was about half past one—three bells in the sea phrase—that the two boats went ashore from the \textit{Hispaniola}. The captain, the squire, and I were talking matters over in the cabin. Had there been a breath of wind, we should have fallen on the six mutineers who were left aboard with us, slipped our cable, and away to sea. But the wind was wanting; and to complete our helplessness, down came Hunter with the news that Jim Hawkins had slipped into a boat and was gone ashore with the rest.

It never occurred to us to doubt Jim Hawkins, but we were alarmed for his safety. With the men in the temper they were in, it seemed an even chance if we should see the lad again. We ran on deck. The pitch was bubbling in the seams; the nasty stench of the place turned me sick; if ever a man smelt fever and dysentery, it was in that abominable anchorage. The six scoundrels were sitting grumbling under a sail in the forecastle; ashore we could see the gigs made fast and a man sitting in each, hard by where the river runs in. One of them was whistling »Lillibullero.«

Waiting was a strain, and it was decided that Hunter and I should go ashore with the jolly-boat in quest of information.

The gigs had leaned to their right, but Hunter and I pulled straight in, in the direction of the stockade upon the chart. The two who were left guarding their boats seemed in a bustle at our appearance; »Lillibullero« stopped off, and I could see the pair discussing what they ought to do. Had they gone and told Silver, all might have turned out differently; but they had their orders, I suppose, and decided to sit quietly where they were and hark back again to »Lillibullero.«

There was a slight bend in the coast, and I steered so as to put it between us; even before we landed we had thus lost sight of the gigs. I jumped out and came as near running as I durst, with a big silk handkerchief under my hat for coolness' sake and a brace of pistols ready primed for safety.

I had not gone a hundred yards when I reached the stockade.

This was how it was: a spring of clear water rose almost at the top of a knoll. Well, on the knoll, and enclosing the spring, they had clapped a stout loghouse fit to hold two score of people on a pinch and loopholed for musketry on either side. All round this they had cleared a wide space, and then the thing was completed by a paling six feet high, without door or opening, too strong to pull down without time and labour and too open to shelter the besiegers. The people in the log-house had them in every way; they stood quiet in shelter and shot the others like partridges. All they wanted was a good watch and food; for, short of a complete surprise, they might have held the place against a regiment.

What particularly took my fancy was the spring. For though we had a good enough place of it in the cabin of the \textit{Hispaniola}, with plenty of arms and ammunition, and things to eat, and excellent wines, there had been one thing overlooked—we had no water. I was thinking this over when there came ringing over the island the cry of a man at the point of death. I was not new to violent death—I have served his Royal Highness the Duke of Cumberland, and got a wound myself at Fontenoy—but I know my pulse went dot and carry one. »Jim Hawkins is gone,« was my first thought.

It is something to have been an old soldier, but more still to have been a doctor. There is no time to dilly-dally in our work. And so now I made up my mind instantly, and with no time lost returned to the shore and jumped on board the jolly-boat.

By good fortune Hunter pulled a good oar. We made the water fly, and the boat was soon alongside and I aboard the schooner.

I found them all shaken, as was natural. The squire was sitting down, as white as a sheet, thinking of the harm he had led us to, the good soul! And one of the six forecastle hands was little better.

»There's a man,« says Captain Smollett, nodding towards him, »new to this work. He came nigh-hand fainting, doctor, when he heard the cry. Another touch of the rudder and that man would join us.«

I told my plan to the captain, and between us we settled on the details of its accomplishment.

We put old Redruth in the gallery between the cabin and the forecastle, with three or four loaded muskets and a mattress for protection. Hunter brought the boat round under the stern-port, and Joyce and I set to work loading her with powder tins, muskets, bags of biscuits, kegs of pork, a cask of cognac, and my invaluable medicine chest.

In the meantime, the squire and the captain stayed on deck, and the latter hailed the coxswain, who was the principal man aboard.

»Mr. Hands,« he said, »here are two of us with a brace of pistols each. If any one of you six make a signal of any description, that man's dead.«

They were a good deal taken aback, and after a little consultation one and all tumbled down the fore companion, thinking no doubt to take us on the rear. But when they saw Redruth waiting for them in the sparred galley, they went about ship at once, and a head popped out again on deck.

»Down, dog!« cries the captain.

And the head popped back again; and we heard no more, for the time, of these six very faint-hearted seamen.

  \begin{figure}[p]
\centering
\includegraphics[width=\linewidth]{chap16downdog}
\caption{»Down, dog!« cries the captain}
\end{figure} 

By this time, tumbling things in as they came, we had the jolly-boat loaded as much as we dared. Joyce and I got out through the stern-port, and we made for shore again as fast as oars could take us.

This second trip fairly aroused the watchers along shore. »Lillibullero« was dropped again; and just before we lost sight of them behind the little point, one of them whipped ashore and disappeared. I had half a mind to change my plan and destroy their boats, but I feared that Silver and the others might be close at hand, and all might very well be lost by trying for too much.

We had soon touched land in the same place as before and set to provision the block house. All three made the first journey, heavily laden, and tossed our stores over the palisade. Then, leaving Joyce to guard them—one man, to be sure, but with half a dozen muskets—Hunter and I returned to the jolly-boat and loaded ourselves once more. So we proceeded without pausing to take breath, till the whole cargo was bestowed, when the two servants took up their position in the block house, and I, with all my power, sculled back to the \textit{Hispaniola}.

That we should have risked a second boat load seems more daring than it really was. They had the advantage of numbers, of course, but we had the advantage of arms. Not one of the men ashore had a musket, and before they could get within range for pistol shooting, we flattered ourselves we should be able to give a good account of a half-dozen at least.

The squire was waiting for me at the stern window, all his faintness gone from him. He caught the painter and made it fast, and we fell to loading the boat for our very lives. Pork, powder, and biscuit was the cargo, with only a musket and a cutlass apiece for the squire and me and Redruth and the captain. The rest of the arms and powder we dropped overboard in two fathoms and a half of water, so that we could see the bright steel shining far below us in the sun, on the clean, sandy bottom.

By this time the tide was beginning to ebb, and the ship was swinging round to her anchor. Voices were heard faintly halloaing in the direction of the two gigs; and though this reassured us for Joyce and Hunter, who were well to the eastward, it warned our party to be off.

Redruth retreated from his place in the gallery and dropped into the boat, which we then brought round to the ship's counter, to be handier for Captain Smollett.

»Now, men,« said he, »do you hear me?«

There was no answer from the forecastle.

»It's to you, Abraham Gray—it's to you I am speaking.«

Still no reply.

»Gray,« resumed Mr Smollett, a little louder, »I am leaving this ship, and I order you to follow your captain. I know you are a good man at bottom, and I dare say not one of the lot of you's as bad as he makes out. I have my watch here in my hand; I give you thirty seconds to join me in.«

There was a pause.

»Come, my fine fellow,« continued the captain; »don't hang so long in stays. I'm risking my life and the lives of these good gentlemen every second.«

There was a sudden scuffle, a sound of blows, and out burst Abraham Gray with a knife cut on the side of the cheek, and came running to the captain like a dog to the whistle.

»I'm with you, sir,« said he.

And the next moment he and the captain had dropped aboard of us, and we had shoved off and given way.

We were clear out of the ship, but not yet ashore in our stockade.